\chapter{Discussione}
\label{ch:discussion}

\section{Interpretazione delle stime}
\label{sec:interpretation}

Da redigere una volta compilato il \Cref{ch:results}. L'interpretazione dovrà
essere disciplinata da due elementi stabiliti in anticipo.

Anzitutto il segno e l'ordine di grandezza attesi in base alla teoria.
\textcite{branger2014} e \textcite{bohringer2012} implicano una risposta dei
volumi misurabile ma contenuta, concentrata nei codici coperti a più alta
intensità emissiva. Un coefficiente assai maggiore segnalerebbe più
probabilmente un controllo contaminato (\Cref{sec:substitution}) che non una
politica insolitamente efficace.

In secondo luogo la possibilità di un risultato nullo. Alla luce di
\textcite{martin2016} e \textcite{naegele2019}, un esito nullo è un'eventualità
concreta e costituirebbe un risultato anziché un fallimento --- in particolare
perché il regime definitivo è il primo contesto in cui l'assegnazione gratuita
viene ritirata invece di schermare la stima. Un risultato nullo, qui, dice
qualcosa che i precedenti esiti nulli sull'\acrshort{ets} non potevano dire.

\section{Meccanismi}
\label{sec:mechanisms}

Tre meccanismi potrebbero generare una stima negativa sui volumi, e i dati
doganali li distinguono solo in parte:

\begin{description}
  \item[Riallocazione effettiva] Gli acquirenti europei passano a origini meno
    intensive o all'offerta interna. Prevede una diminuzione delle quote delle
    origini ad alta intensità \emph{e} un aumento di quelle a bassa intensità.
  \item[\Glsentrytext{resourceshuffling}] Gli esportatori dirottano verso l'UE
    la propria produzione più pulita senza modificarla. Prevede una riduzione
    dell'intensità \emph{dichiarata} senza variazioni nelle quote per origine
    --- e richiede dati di produzione a livello di origine per essere
    confermata, di cui questa tesi non dispone.
  \item[Spostamento a valle] Gli acquirenti si spostano verso beni finiti non
    coperti. Prevede un calo dei codici coperti con un corrispondente aumento a
    valle.
\end{description}

La posizione onesta è che questo disegno può accertare o escludere il terzo
meccanismo in via descrittiva, può distinguere il primo dal secondo soltanto
sotto un'ipotesi di invarianza della produzione delle origini nel breve
periodo, e deve dirlo apertamente anziché presentare la riallocazione come
acquisita.

\section{Implicazioni di policy}
\label{sec:policy}

Il risultato sulla traslazione della DR3 è quello dal contenuto di policy più
diretto. Se il \acrshort{cbam} viene in larga parte assorbito dagli esportatori
extra-UE, esso funziona come un trasferimento di ragioni di scambio a favore
dell'Unione e l'economia politica di una sua estensione è lineare. Se invece
viene in larga parte traslato, si tratta di un costo a carico dei produttori
europei a valle, e diventa assai più forte l'argomento per estenderne l'ambito
verso valle anziché limitarsi ad ampliarlo fra gli input.

\section{Limiti}
\label{sec:limitations}

\begin{enumerate}
  \item \textbf{Periodo post-trattamento breve.} Due o tre trimestri. Nulla di
    quanto qui esposto riguarda la riorganizzazione di lungo periodo delle
    catene di fornitura.
  \item \textbf{Stima relativa, non assoluta.} L'effetto è misurato rispetto ai
    prodotti contigui. Eventuali effetti di anticipazione sull'intero sistema
    economico vengono eliminati dalla differenziazione, sicché la stima
    costituisce un limite inferiore dell'effetto complessivo.
  \item \textbf{Assenza di dati d'impresa.} Comext è a livello
    prodotto--origine--mese. Quali importatori si siano adeguati, e se
    l'adeguamento si sia concentrato fra i grandi dichiaranti, resta qui non
    osservabile.
  \item \textbf{Intensità emissiva imputata.} La $e_{io}$ dell'\Cref{eq:dose}
    proviene da parametri settoriali e non da dati verificati di impianto.
    L'errore di misura nella dose attenua le stime a intensità.
  \item \textbf{Una sola giurisdizione.} La generalizzabilità ad altri
    meccanismi di adeguamento alla frontiera non è verificata e non va
    asserita.
\end{enumerate}

\section{Sviluppi futuri}
\label{sec:further-work}

Il disegno è deliberatamente rieseguibile. Man mano che le fasi successive
estenderanno l'allegato~I ad altri codici \acrshort{cn}, ogni estensione
fornirà una nuova data di trattamento con un gruppo di controllo costruito
secondo la medesima regola --- trasformando un singolo esperimento naturale in
uno scaglionato, punto in cui la strumentazione di \textcite{callaway2021}
diventa necessaria anziché facoltativa. L'integrazione delle dichiarazioni
verificate del registro, una volta pubblicate, sostituirebbe inoltre la
$e_{io}$ imputata con valori osservati e consentirebbe di risolvere, e non
soltanto di delimitare, la questione della riallocazione contabile.
